% ECONOMICS_PROBLEM: {"id": "C61-18", "title": "Polynomial-time solution of rational P-matrix linear complementarity", "jel_code": "C61", "source_id": "OP-0144"}
% !TeX program = xelatex
\documentclass[11pt,a4paper]{article}
\usepackage[margin=23mm]{geometry}
\usepackage{fontspec}
\setmainfont{Latin Modern Roman}
\usepackage{amsmath,amssymb,mathtools}
\usepackage{xcolor,enumitem,booktabs,longtable,array,xurl,hyperref}
\definecolor{muted}{HTML}{53636C}
\hypersetup{colorlinks=true,urlcolor=blue,linkcolor=blue}
\setlength{\parindent}{0pt}
\setlength{\parskip}{5pt}
\newcommand{\R}{\mathbb R}
\newcommand{\E}{\mathbb E}
\newcommand{\N}{\mathbb N}
\newcommand{\ind}{\mathbf 1}
\newcommand{\Var}{\operatorname{Var}}
\newcommand{\Cov}{\operatorname{Cov}}
\newcommand{\supp}{\operatorname{supp}}
\DeclareMathOperator*{\argmin}{arg\,min}
\DeclareMathOperator*{\argmax}{arg\,max}
\newcommand{\entrymeta}[3]{{\sffamily\small\color{muted}\textbf{\texttt{#1}}\quad|\quad #2\par #3\par}}
\newcommand{\Problemref}[1]{\texttt{#1}}
\newcommand{\JELref}[2]{\texttt{#2} (JEL \texttt{#1})}
\begin{document}
\section*{Polynomial-time solution of rational P-matrix linear complementarity}
\textbf{Problem ID:} \texttt{C61-18}\quad \textbf{JEL:} \texttt{C61}\par
% BEGIN ECONOMICS STATEMENT
\label{problem:OP-0144}
% GAME_THEORY_PROBLEM: {"local_id": "GT-014", "original_number": 14, "kind": "P", "entry_kind": "statement_only", "published": false, "review_required": true, "category": "game_theory"}
\label{game:GT-014}
{\sffamily\small\color{muted}
\textbf{GT-014} | Equilibrium-related algebraic search\\
\textbf{P} | Independently corroborated classical promise problem.
\par}
\subsection*{Definitions and mathematical statement}
A matrix $M\in\mathbb Q^{n\times n}$ is a $P$-matrix if every nonempty principal minor is positive:
\[
 \det M[I,I]>0\qquad\text{for all nonempty }I\subseteq\{1,\ldots,n\}.
\]
The input consists of binary-encoded $M$ and $q\in\mathbb Q^n$, with the promise that $M$ is a $P$-matrix. The linear complementarity problem is to find $z,w\in\mathbb Q^n$ with
\[
 w=Mz+q,\qquad z\geq0,\qquad w\geq0,\qquad
 z_iw_i=0\quad(1\leq i\leq n).
\]
Equivalently, the last condition is $z^{\mathsf T}w=0$ because both vectors are nonnegative. The promise guarantees a unique real solution; rational input gives a rational solution. Determine whether an algorithm and a polynomial $p$ exist that solve every promised input using at most $p(L)$ bit operations, where $L$ is the combined input length.
\subsection*{Short English statement}
Can the unique complementarity solution for a rational $P$-matrix always be found in polynomial time?
\subsection*{Variants and scope boundaries}
Checking that an arbitrary matrix is a $P$-matrix is not part of the promise problem. A totalized version may instead output a solution or a witness of promise violation; its output specification differs. Polynomial time in dimension alone, independent of numeric encoding, is a stronger requirement.
\subsection*{Source relationship and status}
Both [S1] and [S2] explicitly retain the polynomial-time question. Their relations to simple stochastic games concern the promised complementarity problem, not recognition of arbitrary $P$-matrices.
\subsection*{References}
\begingroup\footnotesize\raggedright
{\footnotesize\raggedright
\textbf{[S1]} Michaela Borzechowski, John Fearnley, Spencer Gordon, Rahul Savani, Patrick Schnider, and Simon Weber (2024). \emph{Two Choices Are Enough for P-LCPs, USOs, and Colorful Tangents}. ICALP 2024, LIPIcs 297, article 32. Locator: Section 1, ``P-Matrix LCPs''; Definitions 4--7 and Theorem 6. \url{https://drops.dagstuhl.de/entities/document/10.4230/LIPIcs.ICALP.2024.32}\par
\textbf{[S2]} Ioannis Anagnostides and Weiqiang Zheng (2026). \emph{Algorithmic Collusion and the Complexity of Information-Value-Free Equilibria}. arXiv:2609.22757v1, 19 September 2026. Locator: Section 4.2.2, complementarity definition and Theorem 4.2; Section 1.1.2 for the polynomial-time bottleneck. \url{https://arxiv.org/html/2609.22757v1}\par}
\par\endgroup

\subsection*{Common conventions: Source mathematical conventions}

\textbf{Finite games.} For a finite set $A$, $\Delta(A)$ is its probability
simplex. Players choose independent mixed strategies in Nash equilibrium;
correlation is allowed only when explicitly part of the solution concept.
In a game with payoff functions $u_i$ and strategy profile $x$, an additive
$\varepsilon$ Nash equilibrium satisfies
\[
 u_i(x)\geq u_i(x_i',x_{-i})-\varepsilon
 \quad\text{for every player $i$ and every admissible deviation $x_i'$}.
\]
For numerical approximation, payoffs are normalized as locally specified.
An $\varepsilon$ payoff residual is distinct from distance at most
$\varepsilon$ to the set of exact equilibria. Point convergence, convergence
of empirical play, and convergence in distance to an equilibrium set are
also distinct.

\textbf{Computational representations.} Unless stated otherwise, a complexity
question takes rational data with binary encoding; polynomial time is measured
in total bit length. A fixed-accuracy polynomial algorithm is not necessarily
polynomial in $1/\varepsilon$, and neither is necessarily polynomial in
$\log(1/\varepsilon)$. Succinct games and value, demand, or sampling oracles
must declare their interfaces. Exact real solutions require an explicit
representation; an unspecified list of arbitrary real numbers is not a
finite computational output. A claimed algorithmic barrier is meaningful
only for its specified model and reductions.

\textbf{Dynamic games.} Histories, information, observation, and allowable
randomization are part of the model. A stationary strategy depends only on
the current state; a positional strategy in a deterministic graph game
chooses one action at each controlled vertex. Nash equilibrium at an initial
state and equilibrium after every history are different requirements.
An expectation of a pathwise $\liminf$ payoff, a $\liminf$ of expected averages,
a limiting expected average, and a uniform finite-horizon equilibrium are
different criteria. None is silently substituted for another.

\textbf{Allocation and mechanisms.} A complete allocation assigns every item
exactly once unless otherwise stated. Zero-valued goods, negative values,
capacity constraints, feasibility, lotteries, and copies of goods affect
fairness definitions and are specified locally. Dominant-strategy incentive
compatibility, Bayesian incentive compatibility, universal truthfulness,
and truthfulness in expectation impose different inequalities. Whenever
an objective ratio has zero denominator, its convention is stated or those
instances are excluded.

\textbf{Infinite-dimensional models.} Expectations, integrals, stochastic
equations, and optimization objectives require their local regularity,
integrability, filtrations, and admissible domains. A backward--forward
mean-field system is a boundary-value problem; it is not an initial-value
semiflow without an additional construction. Long-time, large-population,
and vanishing-noise limits require an order of limits and a convergence
topology. If a minimum need not be attained, the objective is its infimum
or supremum and the approximation requirement is stated separately.
% END ECONOMICS STATEMENT
\end{document}
